\documentclass[11pt,a4paper]{article}
\usepackage[margin=26mm,headheight=15pt]{geometry}
\usepackage[T1]{fontenc}
\usepackage{amsmath,amsthm,mathtools}
\usepackage{newtxtext,newtxmath}
\usepackage{microtype}
\usepackage{xcolor,booktabs,array,longtable}
\usepackage[most]{tcolorbox}
\usepackage{enumitem}
\usepackage{fancyhdr}
\usepackage{listings}
\usepackage{xurl}
\usepackage[colorlinks=true,linkcolor=navy,citecolor=teal,urlcolor=teal,pdftitle={A Proof of an OEIS Hankel-Numerator Conjecture},pdfauthor={GPT-6 Astra Pro}]{hyperref}
\definecolor{navy}{HTML}{18364A}
\definecolor{teal}{HTML}{216A73}
\definecolor{pale}{HTML}{F0F5F7}
\definecolor{ink}{HTML}{172128}
\color{ink}
\setlength{\parskip}{0.4em}
\setlength{\parindent}{1.2em}
\setlength{\emergencystretch}{2em}
\linespread{1.035}
\pagestyle{fancy}
\fancyhf{}
\fancyhead[L]{\small\scshape Hankel numerators and denominator divisibility}
\fancyhead[R]{\small Research article}
\fancyfoot[L]{\footnotesize 19 September 2026}
\fancyfoot[R]{\footnotesize\thepage}
\renewcommand{\headrulewidth}{0.3pt}
\renewcommand{\footrulewidth}{0pt}
\numberwithin{equation}{section}
\newtheorem{theorem}{Theorem}[section]
\newtheorem{lemma}[theorem]{Lemma}
\newtheorem{proposition}[theorem]{Proposition}
\newtheorem{corollary}[theorem]{Corollary}
\theoremstyle{definition}
\newtheorem{definition}[theorem]{Definition}
\newtheorem{example}[theorem]{Example}
\theoremstyle{remark}
\newtheorem{remark}[theorem]{Remark}
\newcommand{\N}{\mathbb{N}}
\newcommand{\Z}{\mathbb{Z}}
\newcommand{\Q}{\mathbb{Q}}
\newcommand{\vp}{v_p}
\newcommand{\vq}{v_q}
\newcommand{\ord}{\operatorname{ord}}
\newcommand{\num}{\operatorname{num}}
\newcommand{\den}{\operatorname{den}}
\newcommand{\lcm}{\operatorname{lcm}}
\newcommand{\PP}{\Pi}
\newcommand{\floor}[1]{\left\lfloor#1\right\rfloor}
\newcommand{\abs}[1]{\left|#1\right|}
\newcommand{\OEIS}[1]{\href{https://oeis.org/#1}{\texttt{#1}}}
\newtcolorbox{resultbox}[1][]{colback=pale,colframe=teal,boxrule=.5pt,arc=1mm,left=3mm,right=3mm,top=2mm,bottom=2mm,#1}
\lstset{language=Python,basicstyle=\ttfamily\small,keywordstyle=\color{navy}\bfseries,commentstyle=\color{teal},breaklines=true,columns=fullflexible,keepspaces=true,showstringspaces=false,frame=single,rulecolor=\color{teal},xleftmargin=1mm,xrightmargin=1mm}
\setcounter{tocdepth}{1}

\begin{document}
\hypersetup{pageanchor=false}
\begin{titlepage}
\setlength{\parindent}{0pt}
\thispagestyle{empty}
\vspace*{10mm}
{\color{teal}\sffamily\bfseries RESEARCH IN INTEGER SEQUENCES\par}
\vspace{7mm}
{\color{navy}\LARGE\bfseries A Proof of an OEIS\\[2mm] Hankel-Numerator Conjecture\par}
\vspace{5mm}
{\large Prime-power cancellation, sharp denominator divisibility,\\and boxed plane partitions\par}
\vspace{8mm}
{\large GPT-6 Astra Pro\par}
{19 September 2026\par}
\vspace{9mm}
\begin{abstract}
The OEIS entry A122251 records a conjecture of Paul Barry about the reduced numerators of the Hankel determinants of the sequence $1/(pk+1)$, for prime $p$. We prove that formula, and extend it to every positive progression step and every coprime positive offset. For
\[
 H_n(m,a)=\det_{0\leq i,j\leq n}\frac1{a+m(i+j)},\qquad \gcd(a,m)=1,
\]
the reduced numerator is independent of $a$ and equals
\[
 \prod_{p\mid m}p^{\,n(n+1)v_p(m)+2\sum_{k=0}^n v_p(k!)}.
\]
The decisive step is an exact count of residue classes in a square: it proves that every prime not dividing $m$ cancels from the Cauchy determinant numerator. The same count gives all denominator valuations and a sharp greatest-common-divisor theorem over the allowed offsets. A second proof uses an explicit inverse matrix. Further consequences include multiplicative recurrences, digit-sum asymptotics, a Lambert-series generating function with a natural boundary, and a connection with MacMahon's boxed plane partitions. The accompanying standard-library Python implementation passes 123,298 exact checks, including 2,304 independent matrix determinant calculations.
\end{abstract}
\vspace{4mm}
\begin{resultbox}
\small\textbf{Scope of the claim.} This article supplies proofs of the displayed conjecture and its stated extensions. The OEIS entries consulted on 19 September 2026 still label the relevant assertions as conjectural. That is evidence for selecting the problem, not a guarantee of historical priority: an earlier proof or a general result implying it may exist. The classical determinant and inverse-matrix identities are explicitly attributed. No peer review or proof-assistant verification is claimed.
\end{resultbox}
\vfill
{\small\color{teal}Self-contained main proofs \;\textbullet\; Exact arithmetic \;\textbullet\; Reproducible research archive\par}
\end{titlepage}
\hypersetup{pageanchor=true}
\pagenumbering{roman}

\tableofcontents
\vspace{5mm}
\begin{resultbox}
\textbf{Reading guide.} Sections~\ref{sec:problem}--\ref{sec:main} contain a complete elementary proof of the selected conjecture. Section~\ref{sec:denominators} proves the stronger denominator theorem. The remaining sections give independent explanations, sequence-theoretic consequences, examples, and a precise account of the computational checks.
\end{resultbox}
\newpage
\pagenumbering{arabic}

\section{The selected problem and what is proved}\label{sec:problem}

\subsection{A concrete conjecture still marked as such}
For a sequence $(c_k)_{k\geq0}$, its Hankel transform, with the indexing used here, is
\[
 h_n=\det(c_{i+j})_{0\leq i,j\leq n}.
\]
Thus index $n$ means a matrix of size $n+1$. For $c_k=1/(3k+1)$, OEIS entry \OEIS{A122251} describes the resulting numerator sequence as conjectural. Its comments state the more general prime-parameter formula
\begin{equation}\label{eq:original}
 \num\left(\det_{0\leq i,j\leq n}\frac1{1+p(i+j)}\right)
 =(p^2)^{\displaystyle\sum_{k=1}^n\sum_{j=0}^n\floor{k/p^j}}
 \qquad(p\text{ prime}).
\end{equation}
The entry is attributed to Paul Barry, 27 August 2006~\cite{oeis251}. The corresponding $p=2$ entry, \OEIS{A122249}, explicitly requests a proof~\cite{oeis249}. These statements were checked on 19 September 2026.

A targeted search for the sequence identifiers, their formulas, and related generalized Hilbert-matrix arithmetic did not locate a direct prior proof of this exact OEIS assertion. This is not an exhaustive literature survey. In particular, generalized Hilbert and Cauchy matrices have a long literature, and the formulas below are close enough to classical identities that a priority claim would require additional checking.

The reason this is a promising problem is that the determinant itself is readily evaluable. The real issue is \emph{reduction to lowest terms}. A product formula may display many primes in a numerator that disappear completely after cancellation. Our proof isolates and settles precisely that issue.

\subsection{The stronger statement}
Write $v_p(u)$ for the exponent of a prime $p$ in a nonzero integer $u$, extended to nonzero rationals by subtraction. Set
\begin{equation}\label{eq:defH}
 C_n(m,a)=\left(\frac1{a+m(i+j)}\right)_{0\leq i,j\leq n},
 \qquad H_n(m,a)=\det C_n(m,a),
\end{equation}
where $a,m$ are positive integers. Let $A_n(m,a)$ and $B_n(m,a)$ be its positive relatively prime numerator and denominator.

\begin{theorem}[Reduced numerator]\label{thm:numerator}
For $n\geq0$, $a,m\geq1$, and $\gcd(a,m)=1$,
\begin{equation}\label{eq:main}
 \boxed{\displaystyle
 A_n(m,a)=A_n(m)=
 \prod_{p\mid m}p^{\,n(n+1)v_p(m)+2V_p(n)},
 \qquad V_p(n)=\sum_{k=0}^n v_p(k!).}
\end{equation}
The product ranges over the prime divisors of $m$. For $m=1$ it is empty and equals $1$.
\end{theorem}

This theorem proves \eqref{eq:original}, permits composite $m$, and shows that changing the coprime offset does not change the numerator at all. It also shows that the numerator is a square. For $n\geq1$, its prime divisors are exactly those of $m$.

There is a companion statement about denominators. Define
\begin{equation}\label{eq:Ldef}
 L_n=\prod_{k=0}^n(2k+1)\binom{2k}{k}^{\!2}.
\end{equation}
For a positive integer $u$, let $u_{\perp m}$ denote the integer obtained by removing all prime-power factors whose primes divide $m$.

\begin{theorem}[Sharp common denominator divisor]\label{thm:gcd}
For each $n\geq0$ and $m\geq1$,
\begin{equation}\label{eq:gcd}
 \boxed{\displaystyle
 \gcd_{\substack{a\geq1\\\gcd(a,m)=1}} B_n(m,a)=(L_n)_{\perp m}.}
\end{equation}
Moreover, for each prime $q\nmid m$, an explicit offset attains the minimum possible $q$-adic denominator valuation.
\end{theorem}

The infinite gcd means the largest positive integer dividing every denominator in the family. It does \emph{not} assert that a single offset has denominator exactly $(L_n)_{\perp m}$. The primewise minimum is the right notion, and the proof keeps that distinction explicit.

\section{The determinant before cancellation}\label{sec:cauchy}

\subsection{Notation and positivity}
Throughout the proof put
\[
 N=n+1,\qquad P_n=\prod_{k=0}^n k!,\qquad
 Q_n(m,a)=\prod_{i,j=0}^n(a+m(i+j)).
\]
There are only $2n+1$ different factors in $Q_n$, with multiplicities
\begin{equation}\label{eq:Qcompressed}
 Q_n(m,a)=\prod_{s=0}^{2n}(a+ms)^{w_n(s)},
 \qquad w_n(s)=\min(s+1,2n+1-s).
\end{equation}
The determinant is positive. One direct explanation is the Gram representation
\[
 \frac1{a+m(i+j)}=\int_0^1 t^{a-1}t^{mi}t^{mj}\,dt.
\]
For a nonzero real vector $(u_0,\ldots,u_n)$, its quadratic form is the integral of
$t^{a-1}(\sum_i u_it^{mi})^2$. The polynomial is not identically zero, so the integral is positive. Thus the matrix is positive definite and its determinant is positive. Positivity also follows immediately from the product formula below.

\subsection{Cauchy's identity, including a proof}
The classical Cauchy double-alternant identity is a standard determinant evaluation; see Krattenthaler~\cite[Eq.~(2.7)]{krattenthaler}. We give the short algebraic proof to fix signs and indexing.

\begin{lemma}[Cauchy's double alternant]\label{lem:cauchy}
For pairwise distinct $x_i$, pairwise distinct $y_j$, and nonzero $x_i+y_j$,
\begin{equation}\label{eq:cauchy}
 \det_{0\leq i,j<N}\frac1{x_i+y_j}
 =\frac{\displaystyle\prod_{0\leq i<j<N}(x_j-x_i)(y_j-y_i)}
 {\displaystyle\prod_{0\leq i,j<N}(x_i+y_j)}.
\end{equation}
\end{lemma}
\begin{proof}
Multiply the determinant by its displayed common denominator. The result is a polynomial alternating separately in the $x$ variables and in the $y$ variables. It is therefore divisible by both Vandermonde products. Its total degree is $N^2-N$, exactly the sum of their degrees, so it is a constant multiple of their product.

To find the constant, take the residue of the rational identity at $x_{N-1}=-y_{N-1}$. On the determinant side, only the last diagonal entry has this pole; its cofactor is the determinant of size $N-1$. On the product side, the cross factors cancel and leave the same product expression of size $N-1$. Thus the constants for sizes $N$ and $N-1$ agree. At size $1$ the constant is $1$. The identity follows by induction, first for generic variables and then wherever the displayed denominators are nonzero.
\end{proof}

\begin{proposition}[Product evaluation]\label{prop:product}
For all positive $a,m$,
\begin{equation}\label{eq:product}
 H_n(m,a)=\frac{m^{n(n+1)}P_n^2}{Q_n(m,a)}.
\end{equation}
\end{proposition}
\begin{proof}
Take $x_i=mi$ and $y_j=a+mj$ in \eqref{eq:cauchy}. Each Vandermonde is
\[
 \prod_{0\leq i<j\leq n}m(j-i)
 =m^{n(n+1)/2}\prod_{j=0}^n j!.
\]
Multiplying the two Vandermondes gives the numerator in \eqref{eq:product}.
\end{proof}

For later use, the quotient of two consecutive determinants is
\begin{equation}\label{eq:Hratio}
 \frac{H_n(m,a)}{H_{n-1}(m,a)}
 =\frac{m^{2n}(n!)^2}
 {(a+2mn)\displaystyle\prod_{j=0}^{n-1}(a+m(n+j))^2},
\end{equation}
where $H_{-1}=1$ and an empty product is $1$. This follows by separating the last row and column of the denominator product; their common corner is counted only once.

\subsection{Where the difficulty lies}
When $p\mid m$ and $\gcd(a,m)=1$, no factor of $Q_n(m,a)$ is divisible by $p$. All powers of such a prime in the displayed numerator survive. In contrast, if $q\nmid m$, the factor $P_n^2$ can contain a large power of $q$. It is not enough to say that the progression contains multiples of $q$: their \emph{weighted number at every prime-power level} must cover that factorial contribution. The next section supplies an exact count rather than an approximation.

\section{The residue count and the numerator proof}\label{sec:main}

\subsection{A square made of full residue blocks and one corner}
For integers $d\geq1$, $0\leq b<d$, and $0\leq r<d$, define
\begin{equation}\label{eq:Tdef}
 T_{d,b}(r)=\#\{(u,v):0\leq u,v<b,\ u+v\equiv r\pmod d\}.
\end{equation}
Also let
\[
 C_d(N,r)=\#\{(i,j):0\leq i,j<N,\ i+j\equiv r\pmod d\}.
\]

\begin{lemma}[Exact residue-square count]\label{lem:residue}
Write $N=td+b$ with $0\leq b<d$. Then
\begin{align}
 C_d(N,r)&=dt^2+2tb+T_{d,b}(r),\label{eq:rescount}\\
 2\sum_{k=0}^{N-1}\floor{k/d}&=dt(t-1)+2tb,\label{eq:floorsum}\\
 C_d(N,r)-2\sum_{k=0}^{N-1}\floor{k/d}&=dt+T_{d,b}(r)\geq0.\label{eq:surplus}
\end{align}
\end{lemma}
\begin{proof}
Among $0,\ldots,N-1$, each residue modulo $d$ occurs $t$ times, with one extra occurrence for residues $0,\ldots,b-1$. For each residue $u$, the residue paired with it is uniquely $r-u$. The products of the baseline counts contribute $dt^2$. The extra occurrences in the first coordinate contribute $tb$, and those in the second contribute another $tb$. Pairs in which both occurrences are extra contribute $T_{d,b}(r)$. This proves \eqref{eq:rescount}.

For the floor sum, the first $td$ indices split into $t$ blocks of length $d$, on which the floor values are $0,1,\ldots,t-1$. The remaining $b$ indices have floor value $t$. Thus the sum is $dt(t-1)/2+tb$. Subtracting twice this value from \eqref{eq:rescount} gives \eqref{eq:surplus}.
\end{proof}

This identity is the entire cancellation mechanism. The progression offset changes the residue $r$, but the error term is a count and hence cannot become negative.

\subsection{Passing from residue classes to valuations}
We use the elementary factorial valuation identity
\begin{equation}\label{eq:legendre}
 v_q(k!)=\sum_{e\geq1}\floor{k/q^e}.
\end{equation}
Indeed, every multiple of $q$ contributes once, every multiple of $q^2$ contributes one additional time, and so on. Both sides count the powers of $q$ in the factors $1,\ldots,k$.

Fix a prime $q\nmid m$. For each $e\geq1$, let $d=q^e$ and choose
\begin{equation}\label{eq:residueoffset}
 r_e\equiv-a\,m^{-1}\pmod{q^e},\qquad 0\leq r_e<q^e.
\end{equation}
The inverse exists because $q\nmid m$. Exactly $C_{q^e}(N,r_e)$ factors of $Q_n(m,a)$, counted with their original matrix multiplicities, are divisible by $q^e$. Hence
\begin{align}
 v_q(Q_n(m,a))
 &=\sum_{e\geq1}C_{q^e}(N,r_e)\notag\\
 &\geq 2\sum_{e\geq1}\sum_{k=0}^n\floor{k/q^e}
 =2\sum_{k=0}^n v_q(k!)=2V_q(n).\label{eq:cancellation}
\end{align}
All sums are finite: a factor of $Q_n$ is at most $a+2mn$.

\begin{proof}[Proof of Theorem~\ref{thm:numerator}]
If $p\mid m$, coprimality gives $p\nmid a+m(i+j)$ for every pair $(i,j)$. Equation~\eqref{eq:product} therefore yields
\[
 v_p(H_n(m,a))=n(n+1)v_p(m)+2V_p(n).
\]
If $q\nmid m$, equations~\eqref{eq:product} and \eqref{eq:cancellation} give
\[
 v_q(H_n(m,a))=2V_q(n)-v_q(Q_n(m,a))\leq0.
\]
A reduced numerator contains exactly the positive prime valuations of the rational number. The two cases prove \eqref{eq:main}.
\end{proof}

\begin{corollary}[The OEIS conjecture]\label{cor:oeis}
For every prime $p$ and every $n\geq0$, equation~\eqref{eq:original} holds.
\end{corollary}
\begin{proof}
At $m=p$ and $a=1$, Theorem~\ref{thm:numerator} gives exponent
\[
 n(n+1)+2\sum_{k=0}^n\sum_{j\geq1}\floor{k/p^j}
 =2\sum_{k=1}^n\sum_{j\geq0}\floor{k/p^j}.
\]
The $j=0$ terms sum to $n(n+1)/2$. For $n\geq1$ and $j>n$, $p^j>n$, so all further terms vanish. At $n=0$ both sides equal $1$. This is exactly the finite-sum convention in \eqref{eq:original}.
\end{proof}

\subsection{Offsets with a common factor}
The coprimality assumption is transparent rather than restrictive: a common factor can be scaled out.

\begin{corollary}[Arbitrary positive offset and step]\label{cor:noncoprime}
Let $g=\gcd(a,m)$, $a_0=a/g$, and $m_0=m/g$. Then
\begin{equation}\label{eq:noncoprime}
 A_n(m,a)=\prod_{p\mid m_0}
 p^{\max\{0,\ n(n+1)v_p(m_0)+2V_p(n)-(n+1)v_p(g)\}}.
\end{equation}
\end{corollary}
\begin{proof}
Every entry of the matrix is $g^{-1}$ times the corresponding entry of $C_n(m_0,a_0)$. Thus
$H_n(m,a)=g^{-(n+1)}H_n(m_0,a_0)$. Apply Theorem~\ref{thm:numerator} to the coprime pair and subtract $(n+1)v_p(g)$ from its rational valuations. A nonpositive valuation cannot become positive under this scaling. Equivalently, the new numerator is
$A_n(m_0)/\gcd(A_n(m_0),g^{n+1})$.
\end{proof}

\section{Exact denominator valuations and a sharp gcd}\label{sec:denominators}

\subsection{No hidden cancellation remains}
Retaining the exact surplus in Lemma~\ref{lem:residue} gives more information than just the numerator.

\begin{theorem}[Denominator valuation formula]\label{thm:denval}
Assume $\gcd(a,m)=1$ and put $N=n+1$. For a prime $q\mid m$, $v_q(B_n(m,a))=0$. For $q\nmid m$,
\begin{equation}\label{eq:denval}
 v_q(B_n(m,a))=
 \sum_{e\geq1}\left(
 q^e\floor{N/q^e}
 +T_{q^e,N\bmod q^e}(r_e)\right),
\end{equation}
with $r_e$ defined by \eqref{eq:residueoffset}.
\end{theorem}
\begin{proof}
For $q\nmid m$, the proof of Theorem~\ref{thm:numerator} already shows that the rational valuation is nonpositive. The denominator valuation is therefore
$v_q(Q_n)-2V_q(n)$. Apply \eqref{eq:surplus} separately at every $q^e$ and sum. For $q\mid m$ the valuation of the rational determinant is nonnegative, so the reduced denominator has no such prime factor.
\end{proof}

The residual count can be evaluated in constant many arithmetic operations. Set
\[
 h_b(s)=
 \begin{cases}
 s+1,&0\leq s<b,\\
 2b-1-s,&b\leq s\leq2b-2,\\
 0,&\text{otherwise}.
 \end{cases}
\]
For $b=0$ it is identically zero. Since $0\leq u+v\leq2b-2<2d$,
\begin{equation}\label{eq:Ttriangle}
 T_{d,b}(r)=h_b(r)+h_b(r+d).
\end{equation}
The sum in \eqref{eq:denval} may stop when $q^e>a+2mn$. Thus it is both a structural result and an effective algorithm.

\subsection{The least possible residual count}
\begin{lemma}[Sharp residual lower bound]\label{lem:Tmin}
For $0\leq b<d$ and every residue $r$,
\begin{equation}\label{eq:Tmin}
 T_{d,b}(r)\geq\max(0,2b-d),\qquad
 T_{d,b}(d-1)=\max(0,2b-d).
\end{equation}
\end{lemma}
\begin{proof}
Let $I=\{0,\ldots,b-1\}\subseteq\Z/d\Z$. The count is $|I\cap(r-I)|$. Both sets have size $b$, so their intersection has size at least $2b-d$, as well as at least zero.

For $r=d-1$, no sum $u+v$ can equal $2d-1$. Thus only the ordinary equality $u+v=d-1$ is possible. The allowed first coordinates form the interval $d-b\leq u\leq b-1$, whose length is $\max(0,2b-d)$. This attains the bound.
\end{proof}

\subsection{The ordinary Hilbert denominator}
At $m=a=1$, Theorem~\ref{thm:numerator} says the numerator is $1$. Equation~\eqref{eq:Hratio} becomes
\[
 \frac{H_n(1,1)}{H_{n-1}(1,1)}
 =\frac1{(2n+1)\binom{2n}{n}^2}.
\]
Consequently $H_n(1,1)=1/L_n$, with $L_n$ given by \eqref{eq:Ldef}. The first values are
\[
 1,\quad12,\quad2160,\quad6048000,\quad266716800000.
\]
Here the residue $r_e$ equals $q^e-1$ at every level. Theorem~\ref{thm:denval} and Lemma~\ref{lem:Tmin} therefore give
\begin{equation}\label{eq:Lval}
 v_q(L_n)=\sum_{e\geq1}\left[
 q^e\floor{N/q^e}+\max\bigl(0,2(N\bmod q^e)-q^e\bigr)
 \right].
\end{equation}
It follows immediately that $v_q(B_n(m,a))\geq v_q(L_n)$ whenever $q\nmid m$ and $\gcd(a,m)=1$. This proves the divisibility direction of \eqref{eq:gcd}.

\subsection{Constructive sharpness by congruences}
We now show that the lower bound cannot be increased, prime by prime. This is the part needed to obtain an exact gcd rather than just a common divisor.

\begin{proof}[Proof of Theorem~\ref{thm:gcd}]
The preceding subsection proves that $(L_n)_{\perp m}$ divides every allowed denominator. Conversely, fix $q\nmid m$ and choose $K\geq1$ such that
\[
 q^K>2n+1.
\]
The Chinese remainder theorem gives a positive integer $a$ satisfying
\begin{equation}\label{eq:crt}
 a\equiv m\pmod{q^K},\qquad a\equiv1\pmod m.
\end{equation}
In particular $\gcd(a,m)=1$. For $0\leq s\leq2n$,
\[
 a+ms\equiv m(s+1)\pmod{q^K}.
\]
Since $q\nmid m$ and $1\leq s+1<q^K$, the right side has valuation $v_q(s+1)<K$. Congruence modulo $q^K$ then implies
\[
 v_q(a+ms)=v_q(s+1).
\]
The denominator product has exactly the ordinary Hilbert $q$-valuation. Subtracting the same $2V_q(n)$ shows
$v_q(B_n(m,a))=v_q(L_n)$. Thus the minimum over offsets equals $v_q(L_n)$ for every $q\nmid m$. For primes dividing $m$, all denominator valuations are zero. This proves the gcd identity.
\end{proof}

An explicit choice is useful for certificates. Let $M=q^K$ be any prime power exceeding $2n+1$ and choose the least residue
\[
 t\equiv (1-m)M^{-1}\pmod m,\qquad 0\leq t<m.
\]
Then $a=m+Mt$ works. When $m=1$, take $t=0$ and $a=1$. To certify a finite instance of the gcd theorem computationally, start with $B_n(m,1)$ and include one such witness offset for each of its prime divisors. The gcd of these finitely many denominators equals the infinite gcd.

\section{An independent proof from the inverse matrix}\label{sec:inverse}

The residue proof gives the strongest local information, but an inverse-matrix argument explains why the prime support is so restricted. Explicit Cauchy inverses obtained from interpolation are classical; Schechter gives a general treatment~\cite{schechter}. The specialization and its arithmetic application are derived here.

\subsection{A rational binomial integrality lemma}
For rational $x$ and integer $k\geq0$, use the polynomial definition
\[
 \binom{x}{k}=\frac{x(x-1)\cdots(x-k+1)}{k!},\qquad \binom{x}{0}=1.
\]
Say that a rational number is \emph{$q$-integral} if its reduced denominator is not divisible by $q$.

\begin{lemma}\label{lem:qbinomial}
If $q$ is prime and $x=b/c$ with integers $b,c$, $c\neq0$, and $q\nmid c$, then $\binom{x}{k}$ is $q$-integral for every $k\geq0$.
\end{lemma}
\begin{proof}
The assertion is immediate if a numerator factor vanishes. Otherwise,
\[
 \binom{b/c}{k}=\frac{\prod_{j=0}^{k-1}(b-cj)}{c^k k!}.
\]
For each $q^e$, the congruence $b-cj\equiv0\pmod{q^e}$ describes exactly one residue class of $j$. Among $k$ consecutive integers, it occurs at least $\lfloor k/q^e\rfloor$ times. Summing over $e$ proves that the numerator has at least $v_q(k!)$ powers of $q$. The factor $c^k$ is a $q$-unit, so no $q$ remains in the reduced denominator.
\end{proof}

\subsection{The four-binomial inverse}
\begin{proposition}\label{prop:inverse}
Put $\alpha=a/m$. For $0\leq i,j\leq n$,
\begin{align}
 (C_n(m,a)^{-1})_{ij}
 ={}&(-1)^{i+j}m(\alpha+i+j)
 \binom{\alpha+n+i}{n-j}\binom{\alpha+n+j}{n-i}\notag\\[-1mm]
 &\hspace{13mm}\cdot
 \binom{\alpha+i+j-1}{i}\binom{\alpha+i+j-1}{j}.
 \label{eq:inverse}
\end{align}
\end{proposition}
\begin{proof}
First let $G=(1/(\alpha+i+j))_{0\leq i,j\leq n}$, so $C_n=m^{-1}G$. For a fixed column $j$, form the rational interpolation function
\begin{equation}\label{eq:Rj}
 R_j(z)=
 \frac{\prod_{k=0}^n(j+\alpha+k)}{\prod_{\substack{0\leq h\leq n\\h\neq j}}(j-h)}
 \frac{\prod_{\substack{0\leq h\leq n\\h\neq j}}(z-h)}
 {\prod_{k=0}^n(z+\alpha+k)}.
\end{equation}
It satisfies $R_j(h)=\delta_{hj}$ for $h=0,\ldots,n$ and tends to zero at infinity. Its poles are the distinct points $-\alpha-i$. Thus its partial fractions are
\[
 R_j(z)=\sum_{i=0}^n\frac{b_{ij}}{z+\alpha+i}.
\]
Evaluating at each $h=0,\ldots,n$ shows that $Gb_{\bullet j}$ is the $j$th coordinate vector. Hence $b_{ij}=(G^{-1})_{ij}$.

Computing the residue at $z=-\alpha-i$, and using
$\prod_{h\ne j}(j-h)=(-1)^{n-j}j!(n-j)!$, gives
\begin{equation}\label{eq:inverseproduct}
 (G^{-1})_{ij}
 =(-1)^{i+j}
 \frac{\prod_{k=0}^n(\alpha+i+k)\prod_{k=0}^n(\alpha+j+k)}
 {(\alpha+i+j)i!(n-i)!j!(n-j)!}.
\end{equation}
Split each product at its common factor $\alpha+i+j$. For example,
\[
 \binom{\alpha+n+i}{n-j}\binom{\alpha+i+j-1}{j}
 =\frac{\prod_{\substack{0\leq k\leq n\\k\neq j}}(\alpha+i+k)}{(n-j)!j!}.
\]
The analogous identity with $i$ and $j$ interchanged converts \eqref{eq:inverseproduct} into the four-binomial expression. Finally $C_n^{-1}=mG^{-1}$ supplies the factor $m$.
\end{proof}

\begin{corollary}\label{cor:localizedinverse}
Every entry of $C_n(m,a)^{-1}$ lies in $\Z[1/m]$, the rationals whose denominator prime factors all divide $m$. This holds even without $\gcd(a,m)=1$.
\end{corollary}
\begin{proof}
For any prime $q\nmid m$, all upper arguments in \eqref{eq:inverse} are $q$-integral. Lemma~\ref{lem:qbinomial} applies to each binomial. The remaining factor $m(\alpha+i+j)=a+m(i+j)$ is an integer. Hence every entry is $q$-integral for every $q\nmid m$.
\end{proof}

The determinant of this inverse also belongs to $\Z[1/m]$. Therefore
$v_q(H_n(m,a))\leq0$ for all $q\nmid m$. Combined with the exact valuations for $p\mid m$ from Cauchy's product and coprimality, this proves Theorem~\ref{thm:numerator} again. The inverse proof and the residue-square proof share the classical product evaluation for the surviving exponents, but their exclusion of unwanted numerator primes is independent.

\section{A boxed-plane-partition interpretation}\label{sec:partitions}

\subsection{A polynomial in the height}
Define the rational polynomial
\begin{equation}\label{eq:PP}
 \PP_N(t)=\prod_{i=1}^N\prod_{j=1}^N\frac{t+i+j-1}{i+j-1}.
\end{equation}
For a nonnegative integer $t$, MacMahon's boxed-plane-partition formula identifies this product with the number of arrays of size $N\times N$, with entries in $\{0,\ldots,t\}$, weakly decreasing in each row and column. This is the $N\times N\times t$ box specialization of the classical product formula; see the boxed-plane-partition discussion in~\cite{krattenthaler}. We use the known enumeration only in this section, not in the elementary proofs of the main theorems.

\begin{proposition}[Determinant--box identity]\label{prop:box}
For positive $a,m$ and $N=n+1$,
\begin{equation}\label{eq:boxidentity}
 H_n(m,a)^{-1}=m^N L_n\PP_N(a/m-1).
\end{equation}
\end{proposition}
\begin{proof}
Divide the reciprocal of \eqref{eq:product} by its value at $m=a=1$. Since $N^2-n(n+1)=N$, the powers of $m$ leave $m^N$. The remaining product is
\[
 \prod_{i,j=0}^n\frac{a/m+i+j}{1+i+j}
 =\PP_N(a/m-1).
\]
\end{proof}

At $m=1$ and positive integer $a$, Theorem~\ref{thm:numerator} makes the reciprocal determinant equal to the denominator. Hence
\begin{equation}\label{eq:integeroffsetbox}
 B_n(1,a)=L_n\PP_{n+1}(a-1).
\end{equation}
So the denominator divided by the ordinary Hilbert denominator counts boxed plane partitions. At $a=1$ the box height is zero and the quotient is $1$. This gives a combinatorial interpretation of the sharp denominator divisibility in the unscaled case.

\subsection{Why the box polynomial also controls rational offsets}
An integer-valued polynomial need not have integer monomial coefficients. The appropriate basis is the binomial basis.

\begin{lemma}\label{lem:intpoly}
If $P\in\Q[t]$ has degree $D$ and $P(s)\in\Z$ for every nonnegative integer $s$, then
\begin{equation}\label{eq:newton}
 P(t)=\sum_{k=0}^D\Delta^kP(0)\binom{t}{k},
 \qquad \Delta P(t)=P(t+1)-P(t),
\end{equation}
and every coefficient $\Delta^kP(0)$ is an integer. Consequently $P(x)$ is $q$-integral whenever the rational number $x$ is $q$-integral.
\end{lemma}
\begin{proof}
The identity $\Delta\binom{t}{k}=\binom{t}{k-1}$ shows successively that the right side has the same forward differences at zero as $P$, through order $D$. Equivalently, Newton interpolation at $0,\ldots,D$ gives \eqref{eq:newton}; a polynomial of degree at most $D$ vanishing at those $D+1$ points is zero. Each forward difference is an integer linear combination of $P(0),\ldots,P(k)$, so it is integral. Apply Lemma~\ref{lem:qbinomial} term by term for the last assertion.
\end{proof}

The enumerative interpretation makes $\PP_N$ integer-valued on nonnegative integers. If $q\nmid m$, Lemma~\ref{lem:intpoly} therefore gives
$v_q(\PP_N(a/m-1))\geq0$. Equation~\eqref{eq:boxidentity} implies
\[
 v_q(H_n(m,a)^{-1})\geq v_q(L_n).
\]
For coprime $a,m$ this is precisely the universal denominator lower bound already proved by residue counting. Thus an elementary local count, an inverse matrix, and a box-enumeration polynomial illuminate different aspects of the same arithmetic structure.

\section{Recurrences and generating functions}\label{sec:sequences}

\subsection{The prime exponent sequence}
For a prime $p$, define
\begin{equation}\label{eq:Ep}
 E_p(n)=\frac{n(n+1)}2+V_p(n).
\end{equation}
Then $A_n(p)=p^{2E_p(n)}$. The $p=3$ sequence is \OEIS{A122250}; its entry already records the generating function obtained below~\cite{oeis250}. That particular generating function is background, not a novelty claim.

Set $E_p(-1)=E_p(0)=0$. From the definition,
\begin{align}
 E_p(n)-E_p(n-1)&=n+v_p(n!),\label{eq:Epfirst}\\
 E_p(n)-2E_p(n-1)+E_p(n-2)&=1+v_p(n)\qquad(n\geq1).\label{eq:Epsecond}
\end{align}
The $n=1$ case follows with the stated initial convention. Thus the second differences have the particularly simple pattern $1+v_p(n)$.

\begin{center}
\begin{tabular}{r r r r r}
\toprule
$n$ & $E_2(n)$ & $E_3(n)$ & $E_5(n)$ & $E_7(n)$\\
\midrule
0&0&0&0&0\\
1&1&1&1&1\\
2&4&3&3&3\\
3&8&7&6&6\\
4&15&12&10&10\\
5&23&18&16&15\\
\bottomrule
\end{tabular}
\end{center}

\subsection{A multiplicative recurrence for every step}
For positive $r$, define the part supported on primes dividing $m$ by
\[
 r_{(m)}=\prod_{p\mid m}p^{v_p(r)}.
\]
Set $A_{-1}(m)=A_0(m)=1$.

\begin{proposition}\label{prop:multrec}
For every $m\geq1$ and $n\geq1$,
\begin{align}
 \frac{A_n(m)}{A_{n-1}(m)}
 &=m^{2n}\prod_{p\mid m}p^{2v_p(n!)},\label{eq:Aratio}\\
 A_n(m)A_{n-2}(m)
 &=m^2n_{(m)}^2 A_{n-1}(m)^2.\label{eq:multrec}
\end{align}
\end{proposition}
\begin{proof}
Subtract the exponents in \eqref{eq:main} at consecutive indices. The difference of $n(n+1)$ is $2n$, and the difference of $V_p(n)$ is $v_p(n!)$. This proves \eqref{eq:Aratio}. Dividing that identity by its preceding instance leaves $m^2\prod_{p\mid m}p^{2v_p(n)}$, which proves \eqref{eq:multrec}. The initial case $n=1$ is direct.
\end{proof}

In particular $A_{n-1}(m)$ divides $A_n(m)$. For $m>1$, the sequence is strictly log-convex in the multiplicative sense in \eqref{eq:multrec}. At $m=6$, its first factorizations are
\[
 1,\quad 2^2 3^2,\quad 2^8 3^6,\quad 2^{16}3^{14},\quad 2^{30}3^{24}.
\]
The behavior is governed separately by the prime divisors of the step, not by a composite-base floor sum.

\subsection{The Lambert-series and Mahler equations}
Let
\[
 F_p(z)=\sum_{n\geq0}E_p(n)z^n,\qquad
 \mathcal L_p(z)=\sum_{j\geq0}\frac{z^{p^j}}{1-z^{p^j}}.
\]
As formal series, every coefficient of $\mathcal L_p$ is a finite sum. The coefficient of $z^r$, for $r\geq1$, counts the powers $p^j$ dividing $r$, and equals $1+v_p(r)$.

\begin{proposition}\label{prop:gf}
As formal power series, and analytically for $|z|<1$,
\begin{align}
 F_p(z)&=\frac1{(1-z)^2}\sum_{j\geq0}\frac{z^{p^j}}{1-z^{p^j}},\label{eq:gf}\\
 (1-z)^2F_p(z)&=\frac{z}{1-z}+(1-z^p)^2F_p(z^p).\label{eq:mahler}
\end{align}
\end{proposition}
\begin{proof}
Multiply $F_p$ by $(1-z)^2$ and use \eqref{eq:Epsecond}, including the initial coefficients. The result has coefficient $1+v_p(r)$ for $r\geq1$ and zero constant coefficient, hence equals $\mathcal L_p$. Splitting off its $j=0$ term leaves $\mathcal L_p(z^p)$ and gives \eqref{eq:mahler}.

For $|z|\leq\rho<1$, the absolute values of the summands are bounded by $\rho^{p^j}/(1-\rho)$. Their sum converges, giving locally uniform convergence in the open unit disk and justifying the analytic identity there.
\end{proof}

\subsection{A natural boundary and non-holonomy}
\begin{theorem}\label{thm:naturalboundary}
The unit circle is a natural boundary of $F_p$: it has no analytic continuation across any point of that circle. In particular $F_p$ is neither rational nor algebraic, and it satisfies no nontrivial linear differential equation with polynomial coefficients.
\end{theorem}
\begin{proof}
Let $\zeta$ have order $p^K$, with $K\geq1$. Along $z=r\zeta$, $0<r<1$, the terms in $\mathcal L_p$ with $j<K$ have finite limits as $r$ tends to $1$, since $\zeta^{p^j}\neq1$. All terms with $j\geq K$ are positive real numbers,
\[
 \frac{r^{p^j}}{1-r^{p^j}},
\]
and the $j=K$ term alone tends to $+\infty$. Thus the bounded finite head cannot cancel the divergent tail. Since $(1-r\zeta)^2$ tends to a nonzero number, $F_p(r\zeta)$ is unbounded. At $\zeta=1$, positivity gives divergence directly.

Roots of unity whose order is a power of $p$ are dense in the unit circle: the mesh of the $p^K$th roots tends to zero. Every arc therefore contains a point through which $F_p$ cannot be holomorphically continued. An analytic continuation across any point would give a neighborhood, and hence a small arc, on which the continuation is holomorphic, a contradiction.

A rational function has only finitely many finite singularities. An algebraic function has only finitely many possible finite branch points or poles, as follows from the leading coefficient and discriminant of a square-free defining polynomial. Finally, a solution of a linear differential equation with polynomial coefficients can be continued along paths avoiding the finitely many zeros of its leading coefficient. Each possibility is incompatible with the dense set of boundary singularities just exhibited.
\end{proof}

A sequence is \emph{P-recursive} if, from some index onward, it satisfies a nonzero linear recurrence with polynomial coefficients. Coefficient extraction translates such a recurrence into a linear differential equation with polynomial coefficients for its ordinary generating function (an initial polynomial can be removed by differentiating). Thus Theorem~\ref{thm:naturalboundary} shows that $E_p(n)$ is not P-recursive.

The numerator sequence is also not P-recursive for $m>1$, but for a different and simpler reason. Its ordinary generating series has radius zero, so it should not be confused with $F_p$.

\begin{proposition}\label{prop:Anonrec}
For fixed $m>1$, the sequence $A_n(m)$ is not P-recursive, and $\sum_{n\geq0}A_n(m)z^n$ has radius of convergence zero.
\end{proposition}
\begin{proof}
Equation~\eqref{eq:Aratio} gives $A_n/A_{n-1}\geq m^{2n}$, so $A_n\geq m^{n(n+1)}$. The root test gives radius zero.

Suppose $\sum_{j=0}^dP_j(n)A_{n+j}=0$ eventually, where $P_d$ is a nonzero polynomial and $d$ is the largest index with nonzero coefficient. If $d=0$, positivity is already a contradiction. Otherwise divide by $P_d(n)A_{n+d}$ at sufficiently large integers. For $j<d$,
\[
 \frac{A_{n+j}}{A_{n+d}}
 \leq m^{-2\sum_{r=n+j+1}^{n+d}r}
 \leq m^{-2n}.
\]
Each ratio $P_j(n)/P_d(n)$ grows at most polynomially, so every earlier term tends to zero. The last term is $1$, giving the contradiction $1=0$ in the limit.
\end{proof}

\section{Digit sums and asymptotic size}\label{sec:asymptotics}

\subsection{An exact digit formula}
Let $s_p(k)$ be the sum of the base-$p$ digits of $k$, and put
$S_p(n)=\sum_{k=1}^n s_p(k)$. Expanding $k=\sum_{j\geq0}d_jp^j$ in \eqref{eq:legendre} gives
\[
 v_p(k!)=\sum_jd_j(1+p+\cdots+p^{j-1})
 =\frac{k-s_p(k)}{p-1}.
\]
Therefore
\begin{equation}\label{eq:digitE}
 E_p(n)=\frac{p}{2(p-1)}n(n+1)-\frac{S_p(n)}{p-1}.
\end{equation}
This formula separates quadratic growth from digit-dependent fluctuations without any cancellation of large rational determinants.

At a complete digit block, $n=p^\ell-1$, each of the $\ell$ digit positions takes every value $0,\ldots,p-1$ equally often among $0,\ldots,p^\ell-1$. Hence
\[
 S_p(p^\ell-1)=\frac{\ell(p-1)p^\ell}{2},
\]
and
\begin{equation}\label{eq:completeblock}
 E_p(p^\ell-1)=
 \frac{p(p^\ell-1)p^\ell}{2(p-1)}-\frac{\ell p^\ell}{2}.
\end{equation}
The digit formula and complete-block formula are useful independent checks on the floor-sum implementation.

\subsection{An elementary error bound of order $n$}
\begin{lemma}\label{lem:digitsasymp}
For fixed prime $p$ and $n\geq2$,
\begin{equation}\label{eq:digitsasymp}
 S_p(n)=\frac{p-1}{2}\,n\log_p n+O_p(n).
\end{equation}
\end{lemma}
\begin{proof}
At digit position $j$, the digits repeat with period $p^{j+1}$, and over a complete period their mean is $(p-1)/2$. The incomplete terminal period causes an error bounded by a constant depending on $p$ times $p^{j+1}$. Sum over $0\leq j\leq\lfloor\log_p n\rfloor$. The errors form a geometric sum of order $n$. The number of positions is $\log_p n+O(1)$, so their main contribution is $(p-1)(n+1)\log_p n/2+O_p(n)$; replacing $n+1$ by $n$ only adds $O_p(\log n)$.
\end{proof}

Combining \eqref{eq:digitE} and \eqref{eq:digitsasymp} yields
\begin{equation}\label{eq:Easymp}
 E_p(n)=\frac{p}{2(p-1)}n(n+1)-\frac12 n\log_p n+O_p(n).
\end{equation}
For the numerator itself we obtain a uniform-in-offset conclusion.

\begin{corollary}\label{cor:Aasymp}
Fix $m>1$. Let $\omega(m)$ be the number of distinct prime divisors of $m$, and let logarithms without a subscript be natural logarithms. Then, for every positive offset coprime to $m$,
\begin{align}
 \log A_n(m,a)
 ={}&\left(\log m+\sum_{p\mid m}\frac{\log p}{p-1}\right)n(n+1)\notag\\
 &\quad-\omega(m)n\log n+O_m(n).
 \label{eq:Aasymp}
\end{align}
\end{corollary}
\begin{proof}
Take logarithms of \eqref{eq:main} and insert
$V_p(n)=n(n+1)/(2(p-1))-S_p(n)/(p-1)$. Each prime divisor contributes a term $-n\log n$ by \eqref{eq:digitsasymp}. There are finitely many such primes, depending only on $m$, and the numerator is exactly independent of the coprime offset.
\end{proof}

The $O(n)$ remainder can contain digital fluctuations. No smoother expansion or claimed cancellation of those fluctuations is needed for any theorem here.

\section{Examples, tempting false extensions, and limitations}\label{sec:examples}

\subsection{Recovering the conjectured values}
For $p=3$, the exponents begin
\[
 E_3(n)=0,1,3,7,12,18,26,35,45,\ldots.
\]
Thus the numerator sequence begins
\[
 1,\quad9,\quad729,\quad4782969,\quad282429536481,\quad
 150094635296999121,\ldots,
\]
as in \OEIS{A122251}. One exact determinant is
\[
 H_3(3,1)=\frac{4782969}{36653693440000}.
\]
Changing the offset illustrates the stronger theorem:
\begin{center}
\begin{tabular}{r r r}
\toprule
$a$ & $A_3(3,a)$ & $B_3(3,a)$\\
\midrule
1 & 4782969 & 36653693440000\\
2 & 4782969 & 371537512192000\\
4 & 4782969 & 11381495437312000\\
5 & 4782969 & 46222644130432000\\
\bottomrule
\end{tabular}
\end{center}
These are exact rational computations, not decimal approximations.

At $n=3$, $L_3=6048000=2^8 3^3 5^3 7$. The sharp gcds from Theorem~\ref{thm:gcd} are therefore
\begin{center}
\begin{tabular}{r r l}
\toprule
$m$ & $\gcd_a B_3(m,a)$ & Explanation\\
\midrule
1 & 6048000 & No primes are removed.\\
2 & 23625 & Remove the factor $2^8$.\\
3 & 224000 & Remove the factor $3^3$.\\
6 & 875 & Remove both $2^8$ and $3^3$.\\
10 & 189 & Remove both $2^8$ and $5^3$.\\
\bottomrule
\end{tabular}
\end{center}
In every row the gcd ranges only over positive offsets coprime to that row's $m$.

\subsection{Composite substitution is not the right generalization}
Replacing the prime $p$ by an arbitrary composite $m$ in the OEIS floor sum is false. At $m=4$ and $n=2$,
\[
 H_2(4,1)=\frac{16384}{52360425},\qquad A_2(4)=2^{14}.
\]
The naive composite-base formula instead gives
\[
 (4^2)^{\sum_{k=1}^2\sum_{j\geq0}\lfloor k/4^j\rfloor}
 =16^3=4096.
\]
The theorem correctly keeps the prime valuation $v_2(4)=2$ separate from the factorial valuation: the exponent is $2\cdot3\cdot2+2v_2(0!1!2!)=14$.

Coprimality also matters for the square conclusion. With $n=2,m=4,a=2$, scaling out $g=2$ gives
\[
 H_2(4,2)=\frac{32}{496125}.
\]
Its reduced numerator is not a square, exactly as predicted by Corollary~\ref{cor:noncoprime}.

\subsection{Consecutive indices and reciprocal first powers matter}
Arbitrary minors need not share the numerator support of the consecutive Hankel determinants. Selecting indices $\{0,2\}$ from the ordinary Hilbert moment matrix gives
\[
 \det\begin{pmatrix}1&1/3\\1/3&1/5\end{pmatrix}=\frac4{45}.
\]
The underlying moment progression has step $1$, but this sparse selection has introduced an effective spacing of $2$. It is not a counterexample to the theorem; it shows why one cannot omit the consecutive-index condition while leaving $m$ unchanged.

Likewise, squaring each reciprocal moment destroys the Cauchy structure:
\[
 \det\begin{pmatrix}1&1/4\\1/4&1/9\end{pmatrix}=\frac7{144}.
\]
An unexpected numerator prime already appears in size $2$. No theorem about higher reciprocal powers is being claimed.

\subsection{Directions not settled by this article}
For nonconsecutive index sets, the Cauchy product contains the corresponding difference products rather than consecutive factorial products. One can still ask for an exact residue criterion for cancellation, but the uniform square-block argument no longer applies. For moments $(a+mk)^{-r}$ with $r\geq2$, one needs a different determinant evaluation or a different arithmetic invariant. These are possible follow-up questions, not claims that the entire surrounding literature leaves them open.

The present paper settles the precise numerator formula \eqref{eq:original}. It does not identify the minimum numerical denominator among offsets, prove that the gcd is itself attained as one denominator, or classify all related Cauchy minors. It also does not establish priority for any extension merely by failing to find it in a targeted search.

\section{Exact computation and reproducibility}\label{sec:computation}

\subsection{Two independent routes to a rational determinant}
The archive contains \texttt{code/hankel\_arithmetic.py} and \texttt{code/verify.py}. Both use only the Python standard library. Every determinant comparison uses integers and \texttt{fractions.Fraction}; no floating-point tolerance or rounding rule enters the verification.

The direct determinant routine performs Gaussian elimination with row pivoting and rational arithmetic. It does not call the Cauchy product. The product routine evaluates \eqref{eq:product}, using the $2n+1$ different factors in \eqref{eq:Qcompressed}. The theorem routine computes the numerator's prime factorization without constructing the matrix. These three implementations have different failure modes and are compared explicitly.

For a supplied prime $p$, the floor sum is computed without looping over all $k$. If $n+1=td+b$, then
\[
 \sum_{k=0}^n\floor{k/d}=\frac{dt(t-1)}2+tb.
\]
Consequently $V_p(n)$ needs only $O(\log_p n)$ prime-power iterations. Given a factorization of $m$, the factored numerator needs
$O(\sum_{p\mid m}(1+\log_p n))$ arithmetic steps. This counts integer arithmetic operations, not bit operations. The implementation uses trial division to factor the modest parameter $m$; no efficient large-integer factoring claim is intended. Expanding the numerator into decimal form also necessarily costs at least its output size, which grows quadratically in $n$ for fixed $m>1$.

\subsection{What was actually run}
The supplied verification run on Python 3.13.5 passed \textbf{123,298 exact checks}. Its logged elapsed time was 2.448 seconds in the execution environment used for this article; this is a run observation, not a performance guarantee. The important test ranges are as follows.

\begin{center}
\small
\begin{tabular}{@{}p{.31\textwidth}p{.47\textwidth}r@{}}
\toprule
Check family & Parameter range or interpretation & Checks\\
\midrule
Direct matrix vs. product & $0\leq n\leq8$; $1\leq m,a\leq16$ & 2,304\\
Direct matrix vs. theorem & The same independently eliminated matrices & 2,304\\
Extended product vs. theorem & $0\leq n\leq32$; $1\leq m,a\leq24$ & 19,008\\
Universal divisibility & Coprime pairs in the preceding range & 11,847\\
Residue-square formula & $1\leq N\leq30$, $1\leq d\leq40$, every residue & 24,600\\
Residual lower bound & The same range & 24,600\\
Minimum residual at $-1$ & Every $(N,d)$ in that range & 1,200\\
Denominator valuations & $0\leq n\leq10$, $1\leq m,a\leq12$, coprime; primes $\leq47$ & 15,015\\
Inverse identity & Entrywise products for $0\leq n\leq5$, $1\leq m,a\leq8$ & 5,824\\
Inverse denominator support & Every inverse entry in the same range & 5,824\\
Sharpness witnesses & Explicit offsets in finite gcd certificates & 457\\
Finite gcd certificates & $0\leq n\leq6$, $1\leq m\leq12$ & 84\\
\bottomrule
\end{tabular}
\end{center}

Additional checks cover the box-polynomial identity (448), its integer-height integrality (168), the digit formula (3,078), factorial valuation sums (3,078), second differences (3,072), complete digit blocks (30), the multiplicative recurrence (348), and the three boundary counterexamples. For each of the six primes $2,3,5,7,11,13$, a separate Lambert-series coefficient construction is compared with all values through $n=512$; the log counts this as six whole-list checks, not as 3,078 additional checks. All categories and their exact totals are retained in \texttt{data/verification.json}.

These tests establish that the shipped implementation agrees with the formulas over the stated finite ranges. They do not prove an infinite theorem, replace the written proofs, demonstrate independent peer review, or certify historical novelty.

\subsection{Reproduction commands and archive contents}
From the extracted archive directory, run
\begin{lstlisting}[language=bash]
python code/verify.py
python code/hankel_arithmetic.py 3 3 1 --matrix-check
latexmk -pdf -interaction=nonstopmode -halt-on-error report.tex
\end{lstlisting}
The second command prints
\begin{lstlisting}[language={}]
H_3(3,1) = 4782969/36653693440000
numerator factorization: {3: 14}
denominator: 36653693440000
Independent exact matrix check: PASS
\end{lstlisting}
The validation script regenerates the CSV data and logs. An alternate output directory can be chosen with \texttt{--output PATH}. The full source is provided rather than pseudocode alone.

Besides this \LaTeX{} source and compiled PDF, the archive includes numerator and exponent tables, explicit denominator-gcd witnesses, independently computed terms for the two OEIS numerator sequences, a source-status note, and a proposed OEIS comment. None of these files has been submitted to OEIS or any external repository as part of this task.

\section{Conclusion}\label{sec:conclusion}
The conjecture selected from \OEIS{A122251} has a complete elementary proof: Cauchy's product reduces it to a prime-power divisibility question, and a square residue count proves the needed cancellation. The resulting theorem is stronger than the prime-parameter assertion. It covers all progression steps, all coprime offsets, and the noncoprime case after scaling.

The same local count yields exact denominator valuations and the sharp gcd over offsets. An explicit Cauchy inverse gives another explanation of the numerator's restricted prime support. The plane-partition polynomial explains the universal Hilbert denominator factor, while digit sums and Lambert series describe the prime exponents.

The mathematical conclusion is definite within the proofs presented here. The historical conclusion is deliberately narrower: these formulas were still labeled as conjectural in the consulted OEIS entries, and this article supplies a proof, but it does not certify that the proof or the generalizations are the first in the literature. That distinction is especially important when an apparently open database assertion is resolved by combining classical identities with an elementary counting lemma.

\appendix
\section{A compact proof suitable for an OEIS comment}\label{app:comment}
The following is a proposed mathematical comment, not an accepted edit.

Let $H_n=\det_{0\leq i,j\leq n}(1/(1+p(i+j)))$ and $P_n=\prod_{k=0}^n k!$. Cauchy's determinant formula gives
\[
 H_n=\frac{p^{n(n+1)}P_n^2}{\prod_{i,j=0}^n(1+p(i+j))}.
\]
The denominator has no factor $p$, so the surviving $p$-exponent is
$n(n+1)+2\sum_{k=0}^n v_p(k!)$.
For a prime $q\neq p$ and $d=q^e$, set $N=n+1=td+b$ and let $r\equiv-p^{-1}\pmod d$. The number of denominator factors divisible by $d$ is
\[
 dt^2+2tb+\#\{0\leq u,v<b:u+v\equiv r\pmod d\}
 \ \geq\ dt(t-1)+2tb
 =2\sum_{k=0}^n\floor{k/d}.
\]
Sum over $e$ to see that the denominator cancels every $q$-factor of $P_n^2$. The reduced numerator is therefore
\[
 p^{n(n+1)+2\sum_{k=0}^n v_p(k!)}
 =(p^2)^{\sum_{k=1}^n\sum_{j\geq0}\lfloor k/p^j\rfloor},
\]
proving the conjecture for every prime $p$.

\section{A minimal numerator algorithm}\label{app:algorithm}
The implementation in the archive includes input checking, factorization, noncoprime offsets, independent determinants, inverse matrices, and denominator certificates. The core coprime formula can nevertheless be expressed in a few lines once the prime factorization of $m$ is given:
\begin{lstlisting}
def numerator_from_factors(n, factors):
    """factors maps each prime divisor p of m to v_p(m)."""
    answer = 1
    N = n + 1
    for p, power in factors.items():
        V, d = 0, p
        while d <= n:
            t, b = divmod(N, d)
            V += d * t * (t - 1) // 2 + t * b
            d *= p
        answer *= p ** (n * (n + 1) * power + 2 * V)
    return answer
\end{lstlisting}
Here the offset need not appear because Theorem~\ref{thm:numerator} has already proved its irrelevance under coprimality. This algorithm does not estimate or reconstruct a numerator from numerical determinants; it evaluates the proved integer expression directly.

\begin{thebibliography}{9}
\addcontentsline{toc}{section}{References}
\bibitem{oeis251}
P. Barry, \emph{A122251: Numerators of Hankel transform of $1/(3n+1)$ (conjecture)}, The On-Line Encyclopedia of Integer Sequences, entry dated 27 August 2006. Consulted 19 September 2026. \url{https://oeis.org/A122251}.

\bibitem{oeis249}
P. Barry, \emph{A122249: Numerators of Hankel transform of $1/(2n+1)$}, The On-Line Encyclopedia of Integer Sequences, entry dated 27 August 2006, with a 2024 extension. Consulted 19 September 2026. \url{https://oeis.org/A122249}.

\bibitem{oeis250}
P. Barry, \emph{A122250: Partial sums of A004128}, The On-Line Encyclopedia of Integer Sequences, entry dated 27 August 2006. Consulted 19 September 2026. \url{https://oeis.org/A122250}.

\bibitem{krattenthaler}
C. Krattenthaler, \emph{Advanced Determinant Calculus}, S\'eminaire Lotharingien de Combinatoire \textbf{42} (1999), Article B42q, 67 pp. In particular Eq.~(2.7), the Cauchy double alternant, and the discussion of MacMahon's boxed-plane-partition product. \url{https://arxiv.org/abs/math/9902004}.

\bibitem{schechter}
S. Schechter, \emph{On the Inversion of Certain Matrices}, Mathematical Tables and Other Aids to Computation \textbf{13}, no.~66 (1959), 73--77. \url{https://doi.org/10.2307/2001955}.
\end{thebibliography}
\end{document}
